% frontmatter/abstract.tex
% Replace the paragraph below with your abstract. Keep the heading command and \clearpage.

\ndsuprefatoryheading{ABSTRACT}

% Replace this text with your abstract. Conventially, the abstract  summarizes the purpose of the study, the methods or approach, your key findings, and the conclusion or significance of the work.

% Abstracts for doctoral dissertations may be up to 350 words, and abstracts for master’s papers or master’s theses may be up to 150 words. The Abstract page is required in your document and should be page “iii” of the disquisition. When your disquisition is published, it will be available on Google Scholar, Web of Science, ProQuest, NDLTD, and the NDSU Institutional Repository. For dissertations and theses, this abstract should match the abstract in your ETD on ProQuest.

% If you need assistance with this template, please make an appointment with a graduate consultant at the Center for Writers. If you encounter issues while using our templates, would like to request a feature, or have questions about the Publication Review, please contact the Disquisition Coordinator.

Disquisition templates are designed and maintained by the NDSU Center for Writers for your convenience. This sample document also serves as a guide; replace its instructions and placeholders with your own content before submission. Hyperlinks and URLs use black text without underlining. Review the finished document against the formatting requirements and the style manual of your discipline. See \href{https://www.ndsu.edu/cfwriters/disquisitions}{our website} for more information about:

% Use \url{...} for a visible address or \href{destination}{meaningful link text}
% for a descriptive link. Both retain link functionality with compliant styling.

\begin{itemize}
    \item the \href{https://www.ndsu.edu/cfwriters/disquisitions/graduate-school-publication-review-policies}{Graduate School Publication Review} process, including \href{https://www.ndsu.edu/cfwriters/disquisitions/graduate-school-publication-review-policies\#section-1099663}{semester deadlines} for submission, where you can find \href{https://www.ndsu.edu/cfwriters/disquisitions/graduate-school-publication-review-policies\#section-1064488}{help with writing and formatting} your disquisition, how to \href{https://www.ndsu.edu/cfwriters/disquisitions/graduate-school-publication-review-policies\#section-1064489}{incorporate the style manual of your discipline} into this template, and our policies on \href{https://www.ndsu.edu/cfwriters/disquisitions/graduate-school-publication-review-policies\#section-1064491}{copyright} and \href{https://www.ndsu.edu/cfwriters/disquisitions/graduate-school-publication-review-policies\#section-1064492}{co-authored material};
    \item our \href{https://www.ndsu.edu/cfwriters/disquisitions/dqkb}{full guidelines for formatting, construction, and accessibility}, including guides on \href{https://www.ndsu.edu/cfwriters/disquisitions/dqkb/converting-to-pdf}{converting your document to PDF} and creating \href{https://www.ndsu.edu/cfwriters/disquisitions/dqkb/landscape-pages}{landscape-oriented pages};
    \item your \href{https://www.ndsu.edu/cfwriters/disquisitions/pre-submission-requirements}{pre-submission requirements} and \href{https://www.ndsu.edu/cfwriters/disquisitions/submission-checklist}{pre-submission checklist}.
\end{itemize}

If you need assistance with formatting, please make an appointment with a graduate consultant at the \href{https://ndsu.mywconline.com/}{Center for Writers} or \href{https://itservice.ndsu.edu/support/catalog/items/217}{Instructional Design Center}. If you encounter issues while using our templates, would like to request a feature, or have questions about the Publication Review, please \href{mailto:ndsu.disq.coordinator@ndsu.edu}{contact the Disquisition Coordinator}.

Abstracts for doctoral dissertations may be up to 350 words, and abstracts for master's papers or master's theses may be up to 150 words. The Abstract page is required in your document and should be page ``iii'' of the disquisition. When your disquisition is published, this abstract will be included alongside its entry, making it available on Google Scholar, Web of Science, ProQuest, NDLTD, and the NDSU Institutional Repository. For dissertations and theses, this abstract should match the abstract in your ETD on ProQuest.

\clearpage
